\section{Introduction}
Model training costs have been increasing rapidly, and as training data accumulates over time, much of it becomes redundant. This makes data selection a crucial approach to reduce computational burden and enhance model performance~\cite{Settles2009ActiveLL, ds_Toneva2018AnES}. 
Numerous studies have tackled the challenge of selecting informative subsets of training samples, employing diverse criteria to identify the most impactful data for training~\cite{ds_albalak2022survey}. Most existing methods fall into two categories. The first category proposes a measurement metric and selects the top-K samples based on the scores derived from this metric~\cite{ds_ift_Khashabi2020UnifiedQACF, ds_ift_Yang2025DiversitydrivenDS, ds_ift_Zhao2024LongIM}. Some methods further incorporate similarity measures, such as cosine similarity, to ensure diversity among the selected samples~\cite{ds_ift_Yang2025DiversitydrivenDS}. The second category aims to match certain distributions, such as gradients, embeddings, or other indicative features, to better represent the overall dataset~\cite{ds_killamsetty2021grad,ds_bi_Killamsetty2020GLISTERGB,ds_submodular_Mirzasoleiman2019CoresetsFD,ds_ift_cluster_Zhang2024TAGCOSTG}.

However, despite numerous advances in data selection methods, a critical gap remains: most approaches overlook the computational budget as a key factor shaping the selection process. Empirical studies~\cite{hoffmann2022training, kaplan2020scaling} reveal that optimal model performance depends on balancing training data volume, model scale, and available compute budget, highlighting that data selection is inherently linked to computational constraints.
Moreover, recent studies~\citep{ds_ift_Diddee2024ChasingRI,ds_ift_Shen2024RethinkingDS,ds_ift_Xia2024RethinkingDS} reveal that sophisticated selection strategies often fail to consistently outperform random selection across varied experimental settings. \cite{Yin2024ComputeConstrainedDS} further demonstrate that when computational budgets are explicitly factored in, previously effective data selection approaches rarely remain optimal. {\it We thus argue that computational budget must be integral to data selection strategies}, as it determines the appropriate quantity, quality, and distribution of training data, making it a first-order design decision rather than a fixed hyperparameter.

Our argument is further supported by neural network learning dynamics. Models rely on different granularities of knowledge at various training stages, reflected in distinct data subsets. Rahaman et al.~\cite{Rahaman2018OnTS} showed that networks exhibit a spectral bias, learning low-frequency features before higher-frequency ones. Under limited computational budgets, focusing on data rich in low-frequency features can optimize learning, whereas larger budgets allow leveraging diverse, higher-frequency information. We verify these effects empirically in Section~\ref{sec:freq}.

In this paper, we propose a novel Computational Budget-Aware Data Selection (CADS) method, which we naturally formulate within a bilevel optimization framework. In this framework, the inner loop trains the model under the constraints of the computational budget on a selected subset of training data, while the outer loop focuses on optimizing data selection based on the evaluation of the model trained in the inner loop.
Our technical contributions target two main challenges: the expensive estimation of the Hessian matrix for gradients in the outer loop and the computational burden associated with achieving optimality in the inner loop during iterations. To solve the first challenge, we introduce a probabilistic reparameterization strategy and utilize a Hessian-free policy gradient estimator to compute the gradient efficiently. For the second challenge, we reformulate the inner optimization problem as a penalty term in the outer objective function. Notably, we find that estimating the minimum of a one-dimensional loss is sufficient for calculating the gradient, significantly enhancing computational efficiency.

To meet varying data selection demands in practice, we implement CADS in two variants addressing different operational requirements: CADS-E operates at example-level granularity for precise control, while CADS-S assigns weights to source groups for improved scalability. Experiments reveal performance improvements of up to 14.42\% over baselines across vision and language tasks. CADS delivers 3-20× speedup versus conventional bilevel methods, with acceleration scaling proportionally to compute budget size.

\paragraph{Summary of contributions:}
\begin{itemize}
\item We highlight the crucial role of computational budget in data selection, advocating it as a first-order design factor rather than a fixed hyperparameter, thus addressing a key gap by linking data selection to computational constraints.

\item We propose a novel bilevel optimization framework for Computational Budget-Aware Data Selection (CADS), where the inner loop trains the model under budget constraints on a selected data subset, and the outer loop optimizes data selection based on the trained model’s evaluation.

\item To efficiently solve this bilevel problem, we develop techniques including a probabilistic reparameterization for gradient estimation avoiding expensive Hessian computation, and a reformulation of the inner problem as a penalty term in the outer objective. We further show that estimating a one-dimensional loss minimum suffices for gradient calculation, greatly improving computational efficiency.

\item Extensive empirical evaluations demonstrate that our CADS framework achieves superior performance across multiple datasets and training settings, validating the effectiveness of computational budget-aware data selection.
\end{itemize}



